\section{Objections and replies}\label{sec:objections}

We answer the ten objections that bear most directly on the main claim.

\paragraph{1. ``This is just Hume, or Popper.''}
Our first result,
\cref{thm:k1-certify}, is a precise version of the Humean and Popperian point recalled
in \cref{sec:intro}, and we claim no novelty
for it. The paper adds three things the Humean and Popperian traditions do not
supply. First, a usable method need not \emph{verify} a final theory to succeed: it
could simply settle on one and stay there. \Cref{thm:ud-gold} and
\cref{ud:rem-randomized} show that no usable method is guaranteed, even with any fixed
probability, to do that on every computable world compatible with what has been seen.
Any computable implementation of Popper's falsificationism is such a method, and it
inherits the limit. Second, some underdetermination is permanent (objection~3); Popper's
asymmetry between refutation and verification says nothing about rivals that can
never be refuted. Third, the usual escape from underdetermination, preferring the
simplest survivor, depends on the choice of universal machine
(\cref{thm:con-machine}). The Humean point is the first step of the argument, not
its conclusion.

\paragraph{2. ``A method need only work in our world, not in all worlds.''}
Every method
succeeds on some class of worlds and fails on others; we can grant that a method is
reliable in ours. The question is how anyone could \emph{know} that our world lies
in its class. By the prefix form of \cref{thm:ud-gold}, every finite record---the
whole of what anyone has observed, at any time---is consistent with computable worlds
in which a usable deterministic method never settles (for a randomized method, worlds
in which its chance of settling is as small as one likes). And for any method whose
success class meets every record, such as one that identifies the eventually-zero
worlds, the same record is also consistent with worlds in which it succeeds
(\cref{lem:bias}).
The judgment that ours is one of the good worlds is therefore not drawn from the
records. It is the method's bias, restated as a belief about the world. Belot argues that Bayesian
convergence theorems make a similar confidence rationally mandatory \citep{belot2013}.

\paragraph{3. ``Underdetermination is only transient.''}
Each rival in \cref{thm:ud-infinite} disagrees with the others somewhere, so each is
refuted once the record reaches that setting. Sklar calls this ``transient''
underdetermination \citep[p.~380]{sklar1975}, and the objection is that it is only a temporary
state of the evidence;
Laudan and Leplin add that empirical equivalence is relative to a time and to
auxiliary assumptions, and may later be broken \citep{laudan-leplin1991}. We grant
both points rival by rival. Two things remain. First, at every stage infinitely
many unrefuted rivals still disagree about what comes next, and by the prefix form
of \cref{thm:ud-gold} no usable method can tell when the rivals that matter are
gone: the transient phase never announces its end. Stanford argues that recurring
transient underdetermination, in the form of alternatives no one has thought of, is
the real threat to realism \citep{stanford2006}; the rivals of \cref{thm:ud-infinite} show that unrefuted alternatives always exist, though they
are not Stanford's unconceived theories. Second, some underdetermination is not transient at all. Models that
agree on every possible observation are never separated, and chance hypotheses that
agree on which events have probability zero are never told apart with certainty by
any method (\cref{thm:t5-permanent}).

\paragraph{4. ``The rivals are cheap.''}
Kukla gave algorithms that turn any theory into an empirically equivalent rival
(\cref{sec:certify}) \citep{kukla1993}. Leplin
and Laudan reply that such constructions are not genuine scientific rivals
\citep{leplin-laudan1993}; Kukla later defended them \citep{kukla1996}; and Stanford that rivals produced by a general recipe
amount to Cartesian skepticism \citep{stanford2001}. Three replies. First, the
criteria that exclude cheap rivals---independence, naturalness, explanatory
depth---are exactly the kind of preference \cref{lem:bias} identifies: real, often
useful, and not supplied by the records. Second, the rivals need not be cheap.
\Cref{thm:con-nearmin} gives rivals within $\K(x)+O(1)$ bits of the shortest model that
predict a different outcome at any unobserved setting, so the excluding criterion
cannot be length. Third, physics supplies rivals no one regards as artificial:
Newtonian, Lagrangian and Hamiltonian mechanics; matrix and wave mechanics; dual
descriptions of the same physics; and the interpretations of quantum mechanics
(\cref{sec:simplicity,sec:ideal}). Nor is the argument skeptical: it posits no
deceiver and holds for any record, however it was produced.

\paragraph{5. ``Simplicity, or beauty, tracks the truth.''}
The hope for a final theory is usually a hope that the true theory is simple and
beautiful \citep{weinberg1992}. Five points.
\begin{itemize}
  \item \emph{Description length.} The most precise notion of simplicity depends on
        the choice of universal machine (\cref{thm:con-machine}): on any finite record,
        any consistent model can be made the shortest, at a cost of about $\K(e)$
        bits, so the record leaves open the order of candidates within that margin.
        Description length is called objective because it is invariant up to an
        additive constant. We agree, both asymptotically and for comparisons separated
        by more than that constant; but near-tied candidates that disagree about the
        unobserved always exist, and a change of machine that moves no description
        length by more than about the gap makes any of them the shortest
        (\cref{thm:con-nearmin,thm:con-gap}).
  \item \emph{Ideal MDL.} \citet{vitanyi-li2000} show that ideal MDL is Bayes-optimal
        when the data are random for the true hypothesis. We accept this. It is a
        conditional guarantee whose condition, typicality, no finite record can confirm,
        since randomness deficiency is uncomputable; and it yields convergence in the
        limit, never certification.
  \item \emph{Hard to vary.} Deutsch holds that good explanations are those whose
        details cannot be varied without spoiling their fit, and that choosing them is
        rational, not a matter of taste \citep[ch.~1]{deutsch2011}. Taken informally, this is
        one more constraint in the sense of \cref{lem:bias}; formalized as a count of
        free parameters or a description length, it inherits machine-relativity. But
        Deutsch also denies that any explanation is final: ``there can be no such thing
        as an ultimate explanation'' \citep[ch.~3]{deutsch2011}. His view does not oppose the main
        claim. He opposes us only on whether the selection is rational,
        and nothing in our argument requires that it be irrational, only that the
        records do not make it.
  \item \emph{Ockham efficiency.} Kelly and Genin prove that an Ockham method, which
        prefers the simplest answer to a fixed question, converges to the truth with
        worst-case optimal retractions \citep{kelly2007}, or optimal reversals and
        cycles of opinion \citep{genin-kelly2019}. This
        justification is not machine-relative, and we accept it. It is relative to a
        chosen question and its partition into answers, and it guarantees efficient
        convergence in the limit, never certification: an Ockham method never knows
        that it has stopped retracting. It is consistent with our claim.
  \item \emph{Beauty and naturalness.} Our point is independent of Hossenfelder's
        diagnosis (\cref{sec:intro}):
        such criteria are constraints the records do not supply, and no
        record can confirm that the world satisfies them.
\end{itemize}

\paragraph{6. ``The success of physics shows convergence on the truth.''}
The no-miracles argument holds that the predictive success of mature theories would
be a miracle unless they were approximately true \citep{putnam1975}. We accept the
success and explain it differently. A compact model of what is known commits itself
about what is not, and when those commitments survive test the model has earned its
place (\cref{sec:predict}). The ideal predictor of \cref{thm:pred-relative} enters only as a
benchmark: it shows that weighting models by compressed length is a
bias that loses little against any computable model, but no scientist uses it, and
it does not explain why particular theories succeed. The
success of physics is evidence that its theories are good compressions; it cannot
certify that any of them is the last (\cref{thm:k1-certify,sec:predict}).

\paragraph{7. ``Algorithmic probability is itself the law.''}
Müller's algorithmic idealism \citep{muller2026} shares our starting point and our tools: it begins from
the observer's record rather than from a world, and it is formalized within
algorithmic information theory \citep{muller2020}. It also anticipates \cref{thm:con-machine}. In
\emph{Law without law} \citep{muller2020} Müller holds, citing his own result that the most natural attempt to
privilege a universal machine fails \citep{muller2010}, that the choice of algorithmic prior is
arbitrary, and he restricts the theory's predictions to those that are identical for
every choice. From these postulates he derives that observers will, with high
probability, come to see an external world governed by simple computable laws. The
step we decline is his postulate that algorithmic probability \emph{gives the chances}
of what an observer sees next, which makes it the fundamental law. Read as a candidate
final theory, this postulate is subject to three limits. No
record certifies it: every finite record has positive probability under infinitely
many computable measures that disagree about what comes next
(\cref{thm:ud-infinite}(c)), so by part (b) of \cref{thm:k1-certify} any certification
of the universal mixture at a record is wrong in some computable world that produces
that record.
Adopting it is a choice of bias (\cref{lem:bias}), and Müller's own restriction to
machine-invariant predictions concedes that single-step chances are not fixed by the
theory; its machine-invariant predictions are asymptotic. And the mixture is uncomputable, so no usable agent can apply the law
it would be (\cref{pred:prop:no-computable}). Algorithmic idealism is therefore not a
counterexample to our claim but an instance of it: a constraint,
explicitly chosen, that no observation could confirm as final. The two views agree on
what records can and cannot settle, and differ only on whether to affirm a law beyond
them. Müller affirms one; we remain agnostic.

\paragraph{8. ``The second premise begs the question.''}
Refuting \cref{prem:effective} requires a \emph{usable} hypercomputer: a system whose
non-computable behavior can be specified, relied on and put to work, not one whose
user cannot rule out spurious signals or device failure \citep{earman-norton1993}.
For deterministic methods C2 survives any hypercomputer
with the power of a Turing oracle, provided theories can use the same power, because
the diagonal argument reruns at every level of computing power
(\cref{thm:ud-gold-rel,rem:p2star}). This also answers Faizal and colleagues (\cref{sec:intro}): non-algorithmic
understanding would help only if it were
unavailable to the theories it was used to choose among.

\paragraph{9. ``Knowledge needs only reliability, not proof.''}
On fallibilist and externalist accounts, a belief can be knowledge because it was
reliably produced or is safe, without the record ruling out every rival. The
externalist may say that a reliable method that has settled on the final theory, in a
world where it is final, yields knowledge in their sense. We need not deny this: a
believer in that world may hold the truth. What no one can do is establish, from the
records, that they are in that world (C1, C2, \cref{lem:bias}), and that is what
``known to be final'' means here (\cref{def:known}).

\paragraph{10. ``The argument undermines itself.''}
\Cref{prem:record,prem:effective} are universal empirical claims, so, as for any
universal claim, no record certifies them either. We do not claim to know them. We
assert them as premises and state what would refute each (\cref{sec:framework,rem:p2star};
objection~8), so every conclusion is conditional on them. This is consistent: the paper claims no certification of a kind it
shows to be impossible.
